\documentclass[12pt,a4paper,fontset=none]{ctexart}

\usepackage[a4paper,margin=2.4cm]{geometry}
\usepackage{amsmath}
\usepackage{amssymb}
\usepackage{booktabs}
\usepackage{array}
\usepackage{longtable}
\usepackage{enumitem}
\usepackage{hyperref}
\usepackage{xcolor}
\usepackage{listings}

\setCJKmainfont{Microsoft JhengHei}
\setCJKsansfont{Microsoft JhengHei}
\setCJKmonofont{Microsoft JhengHei}

\hypersetup{
  colorlinks=true,
  linkcolor=blue!55!black,
  urlcolor=blue!55!black,
  citecolor=blue!55!black
}

\lstdefinestyle{cmd}{
  basicstyle=\ttfamily\small,
  breaklines=true,
  frame=single,
  columns=fullflexible
}

\setlist[itemize]{leftmargin=2em}
\setlist[enumerate]{leftmargin=2em}

\title{\textbf{shogiAI 專案報告}\\
\large 將棋 CSA 棋譜資料庫與數據科學模型設計}
\author{專案名稱：shogiAI}
\date{2026 年 5 月}

\begin{document}
\maketitle

\begin{abstract}
本專案以日本將棋 CSA 棋譜為核心資料，建立一套從棋譜蒐集、格式解析、資料庫正規化、局面重建、訓練資料匯出、policy/value 神經網路訓練，到網頁棋譜瀏覽與 AI 對弈的完整流程。系統目前包含 1000 份 CSA 棋譜，並由資料庫匯出 111238 筆局面樣本，每筆樣本包含 $43 \times 9 \times 9$ 的局面特徵、13689 類走法標籤與勝負價值標籤。本報告分成三個層次：第一部分說明整體主系統與製作流程；第二部分以資料庫課程角度說明關聯式設計、正規化、匯入與查詢；第三部分以數據科學課程角度說明特徵工程、監督式學習、CNN policy/value model，以及與 alpha-beta 搜尋整合的理論方法。
\end{abstract}

\tableofcontents
\newpage

\section{主要系統設計}

\subsection{專案目標}
shogiAI 的核心目標不是只製作一個能讀取棋譜的工具，而是建立一條可重複執行的資料管線，使原始 CSA 棋譜能同時服務於資料庫管理、棋譜瀏覽、數據分析與 AI 模型訓練。整體設計可分為下列目標：

\begin{itemize}
  \item 將原始 CSA 棋譜轉換為乾淨、可查詢、可重建局面的資料。
  \item 以 MySQL 保存棋局、棋手、每手走法與每個局面，使棋譜不只是檔案，而是可被條件查詢的結構化資料。
  \item 由資料庫重建訓練資料，產生 policy/value network 可使用的 NumPy 壓縮資料集。
  \item 使用 CNN 學習棋手在特定局面下的走法選擇與勝負估計。
  \item 在網頁介面中提供棋譜瀏覽、自行對弈、AI 對弈、policy 候選手與搜尋資訊。
\end{itemize}

\subsection{專案檔案與模組分工}

\begin{longtable}{p{0.28\textwidth}p{0.62\textwidth}}
\toprule
\textbf{檔案或目錄} & \textbf{功能}\\
\midrule
\endhead
\texttt{data/*.csa} & 原始 CSA 棋譜資料，共 1000 份。\\
\texttt{src/csa\_preprocess.py} & 解析 CSA/KIF、驗證合法走法、編碼局面與走法、輸出 \texttt{.npz} 訓練樣本。\\
\texttt{src/import\_csa\_to\_db.py} & 將 CSA 棋譜匯入 MySQL，建立 \texttt{users}、\texttt{players}、\texttt{game\_records}、\texttt{moves}、\texttt{positions} 五張表。\\
\texttt{src/export\_policy\_dataset.py} & 從 MySQL 讀取局面與走法，重建 policy/value 訓練資料。\\
\texttt{src/policy\_model.py} & 定義 CNN policy/value network 與推論封裝。\\
\texttt{src/train\_policy.py} & 執行資料切分、模型訓練、驗證、測試與 checkpoint 儲存。\\
\texttt{src/ai\_search.py} & 實作將棋局面評估、negamax alpha-beta 搜尋、quiescence search、transposition table、killer/history move ordering。\\
\texttt{src/csa\_browser\_api.py} & 提供 HTTP API，支援從 MySQL 或本機 CSA 檔讀取棋局，並連接 AI 搜尋與 policy model。\\
\texttt{web/index.html}, \texttt{web/app.js}, \texttt{web/styles.css} & 前端介面，包含棋譜瀏覽、自行對弈、AI 對弈與資料庫統計。\\
\texttt{out/policy\_dataset.npz} & 從資料庫匯出的訓練資料，共 111238 筆樣本。\\
\texttt{out/policy\_model.pt} & 已訓練完成的 PyTorch 模型 checkpoint。\\
\bottomrule
\end{longtable}

\subsection{整體資料流程}
整體流程可視為一條 ETL 與機器學習管線：

\begin{enumerate}
  \item \textbf{Extract：} 從 \texttt{data} 目錄讀取 CSA 棋譜。
  \item \textbf{Transform：} 使用 \texttt{cshogi} 解析棋譜，檢查每手是否合法，並在每手之前保存局面。
  \item \textbf{Load：} 將棋局、棋手、走法、局面寫入 MySQL 正規化資料表。
  \item \textbf{Dataset Export：} 從資料庫重新組合局面與下一手，轉換成 CNN 可訓練的張量資料。
  \item \textbf{Model Training：} 訓練 policy/value network，學習走法分類與局面勝率估計。
  \item \textbf{Serving：} 透過 \texttt{csa\_browser\_api.py} 提供棋譜瀏覽、AI 對弈與模型候選手分析。
\end{enumerate}

\subsection{系統架構}
系統架構的設計重點是分離「資料保存」、「模型訓練」與「使用者互動」。CSA 原始檔只作為初始資料來源；MySQL 成為查詢、重建與後續分析的主要資料層；模型與搜尋則透過 API 被前端呼叫。

\[
\text{CSA Files}
\rightarrow
\text{Parser / Validator}
\rightarrow
\text{MySQL}
\rightarrow
\text{NPZ Dataset}
\rightarrow
\text{CNN Model}
\rightarrow
\text{Search + Web API}
\rightarrow
\text{Browser UI}
\]

此架構的優點是每一層都可以獨立測試。若棋譜匯入失敗，可檢查資料庫層；若模型效果不佳，可重新匯出資料或調整訓練；若前端呈現錯誤，可直接呼叫 API 驗證 JSON 回應。

\subsection{製作流程回顧}
本專案的製作流程大致經過下列階段：

\begin{enumerate}
  \item 建立 CSA 棋譜前處理程式，先確認可以把棋譜轉成合法局面與下一手。
  \item 擴充 MySQL 匯入流程，將原本檔案式資料改成資料表設計。
  \item 建立棋譜瀏覽 API 與網頁前端，用資料庫資料還原棋盤、持駒、手數與走法列表。
  \item 由資料庫匯出 policy/value 訓練資料，讓資料科學流程不直接依賴原始檔案。
  \item 建立 CNN 模型，加入 value mask、AdamW、cosine learning rate scheduler、依棋局切分資料集等訓練設計。
  \item 建立 AI 搜尋，將傳統 alpha-beta 搜尋與神經網路 policy ordering、value head 評估結合。
\end{enumerate}

\newpage
\section{資料庫課程報告：CSA 棋譜資料庫設計}

\subsection{資料庫設計目標}
資料庫部分的核心問題是：CSA 棋譜原本是半結構化文字檔，雖然適合保存對局紀錄，但不適合查詢、統計與局面重建。因此資料庫設計需要達成下列目標：

\begin{itemize}
  \item 保存棋局層級資訊，例如賽事、場地、日期、開局、勝負結果與原始檔名。
  \item 保存棋手資訊，避免同一棋手在多個棋局中重複出現文字資料。
  \item 保存每一手走法，使系統能依手數回放棋局。
  \item 保存每一手之後的 SFEN 局面，使前端與資料科學流程可以快速取得任意局面。
  \item 支援條件查詢，例如依賽事、日期區間、棋手、開局搜尋棋局。
  \item 支援去重匯入，避免同一份 CSA 被重複寫入。
\end{itemize}

\subsection{關聯式資料模型}
本專案資料庫共使用五張主要資料表：

\begin{longtable}{p{0.22\textwidth}p{0.68\textwidth}}
\toprule
\textbf{資料表} & \textbf{設計目的}\\
\midrule
\endhead
\texttt{users} & 保存上傳者或匯入者資訊。雖然目前主要由系統匯入，但保留使用者表可支援未來多人上傳。\\
\texttt{players} & 保存棋手姓名與段位。棋手被抽離成獨立資料表，避免在每筆棋局中重複保存字串。\\
\texttt{game\_records} & 保存棋局主檔，包含黑方、白方、賽事、場地、開局、勝負、原始檔名與日期。\\
\texttt{moves} & 保存每一手走法，包含手數、輪到哪一方、CSA/USI 表示、用時與註解。\\
\texttt{positions} & 保存每一手後的局面，包含手數、輪到哪一方、SFEN、hash、是否被將軍與合法手數量。\\
\bottomrule
\end{longtable}

\subsection{ER 關係說明}
主要關聯如下：

\begin{itemize}
  \item \texttt{users.user\_id} 一對多連到 \texttt{game\_records.uploader\_id}。
  \item \texttt{players.player\_id} 分別連到 \texttt{game\_records.black\_player\_id} 與 \texttt{game\_records.white\_player\_id}。
  \item \texttt{game\_records.game\_id} 一對多連到 \texttt{moves.game\_id}。
  \item \texttt{game\_records.game\_id} 一對多連到 \texttt{positions.game\_id}。
\end{itemize}

以概念式 ER 模型表示：

\[
\text{User}
\xrightarrow{1:N}
\text{GameRecord}
\xleftarrow{N:1}
\text{Player}
\]

\[
\text{GameRecord}
\xrightarrow{1:N}
\text{Move},
\qquad
\text{GameRecord}
\xrightarrow{1:N}
\text{Position}
\]

\subsection{正規化設計}
本資料庫設計主要符合第三正規化（Third Normal Form, 3NF）的精神：

\begin{itemize}
  \item 第一正規化：每個欄位保存單一語意，例如棋手姓名、日期、走法、SFEN 分開保存。
  \item 第二正規化：每張表的非鍵欄位都依賴該表主鍵，例如 \texttt{moves} 的 \texttt{usi\_move} 依賴 \texttt{move\_id}。
  \item 第三正規化：棋手姓名不直接重複放在 \texttt{game\_records} 中，而是透過 \texttt{players} 外鍵參照，避免傳遞相依。
\end{itemize}

這樣設計的好處是資料一致性較高。例如同一棋手出現在多場比賽時，資料庫只需保存一份棋手資料。若未來要統計某位棋手的勝負或常用開局，也能直接由 \texttt{players} 與 \texttt{game\_records} 關聯查詢。

\subsection{資料表欄位設計}

\subsubsection{棋局主檔：\texttt{game\_records}}
\texttt{game\_records} 是整個資料庫的中心表。其欄位包含：

\begin{itemize}
  \item \texttt{game\_id}：主鍵。
  \item \texttt{black\_player\_id}, \texttt{white\_player\_id}：黑方與白方棋手外鍵。
  \item \texttt{event\_name}, \texttt{site}, \texttt{opening}：CSA metadata。
  \item \texttt{result}：勝負結果，例如 \texttt{TORYO}, \texttt{DRAW}, \texttt{SENNICHITE}。
  \item \texttt{source\_format}：資料來源格式，目前為 CSA。
  \item \texttt{original\_file\_name}：原始檔名，用於追蹤來源與去重。
  \item \texttt{played\_at}：對局日期。
\end{itemize}

\subsubsection{走法表：\texttt{moves}}
\texttt{moves} 保存棋局的時間序列。每一列是一手棋：

\begin{itemize}
  \item \texttt{move\_number}：第幾手。
  \item \texttt{side\_to\_move}：此手由黑方或白方下。
  \item \texttt{original\_move} 與 \texttt{usi\_move}：保存原始走法與標準化 USI 表示。
  \item \texttt{move\_time}, \texttt{comment}：保存棋譜附加資訊。
\end{itemize}

\subsubsection{局面表：\texttt{positions}}
\texttt{positions} 是本專案資料庫設計中最重要的表之一。傳統棋譜只保存走法，但本專案額外保存每一手後的 SFEN 局面，使系統可以快速取得第 $n$ 手局面，而不必每次從第一手重播。

\begin{itemize}
  \item \texttt{move\_number}：此局面對應到第幾手後；0 代表初始局面。
  \item \texttt{sfen}：完整局面表示，可被 \texttt{cshogi.Board} 直接重建。
  \item \texttt{sfen\_hash}：局面 hash，可用於重複局面偵測與索引。
  \item \texttt{is\_check}：是否正在被將軍。
  \item \texttt{legal\_moves\_count}：該局面合法手數量，可用於分析複雜度。
\end{itemize}

\subsection{匯入流程與交易控制}
\texttt{import\_csa\_to\_db.py} 的匯入流程如下：

\begin{enumerate}
  \item 掃描輸入路徑中的 CSA 檔案。
  \item 使用 \texttt{cshogi.CSA.Parser} 解析棋譜。
  \item 讀取棋手、賽事、日期與結果等 metadata。
  \item 透過 \texttt{ensure\_player} 避免重複建立棋手。
  \item 寫入 \texttt{game\_records} 後取得 \texttt{game\_id}。
  \item 從初始局面開始逐手檢查合法性，建立 \texttt{moves} 與 \texttt{positions}。
  \item 使用交易 \texttt{commit} 確保每局資料完整寫入；若發生錯誤則 rollback。
\end{enumerate}

此設計使用 \texttt{--skip-existing} 依 \texttt{original\_file\_name} 跳過已匯入棋局，並提供 \texttt{--max-games} 與 \texttt{--pause-every} 來因應學校 MySQL 可能存在的查詢限制。

\subsection{查詢與 API 設計}
前端使用 \texttt{csa\_browser\_api.py} 提供的 API 存取資料。主要端點包括：

\begin{longtable}{p{0.34\textwidth}p{0.56\textwidth}}
\toprule
\textbf{API} & \textbf{用途}\\
\midrule
\endhead
\texttt{GET /api/games} & 取得棋局清單，可依賽事、日期、棋手、開局篩選。\\
\texttt{GET /api/games/<id>?ply=N} & 取得指定棋局第 $N$ 手後的局面。\\
\texttt{GET /api/db/stats} & 回傳資料庫棋局數、棋手數、走法數、局面數與重複檔案統計。\\
\texttt{POST /api/self-play/state} & 由前端傳入 USI 走法序列，後端重建自行對弈局面。\\
\texttt{POST /api/ai-play/move} & 呼叫 AI 搜尋取得下一手。\\
\texttt{POST /api/policy/candidates} & 回傳 policy model 對目前合法手的候選排序。\\
\bottomrule
\end{longtable}

在 SQL 查詢上，棋局清單使用 \texttt{LEFT JOIN players} 取得黑白棋手，並使用 \texttt{COUNT(m.move\_id)} 統計手數。篩選條件採用參數化查詢，例如棋手搜尋會對黑方與白方姓名做 \texttt{LIKE} 比對，避免直接拼接使用者輸入造成 SQL injection。

\subsection{資料庫設計評估}
本資料庫設計的優點如下：

\begin{itemize}
  \item \textbf{可重建性：} 保存每一手局面，使任意手數可直接回放。
  \item \textbf{可查詢性：} 棋局、棋手、走法、局面分表後，可針對不同層級做查詢。
  \item \textbf{資料一致性：} 使用外鍵維持棋局與棋手、走法、局面的關係。
  \item \textbf{可擴充性：} 未來可新增 openings、tournaments、model\_predictions 等表。
  \item \textbf{可服務數據科學：} \texttt{positions} 與 \texttt{moves} 可直接 JOIN 成監督式學習資料。
\end{itemize}

後續若要優化資料庫效能，建議加入下列索引：

\begin{lstlisting}[style=cmd]
CREATE INDEX idx_game_records_played_at ON game_records(played_at);
CREATE INDEX idx_game_records_original_file ON game_records(original_file_name);
CREATE INDEX idx_moves_game_number ON moves(game_id, move_number);
CREATE INDEX idx_positions_game_number ON positions(game_id, move_number);
CREATE INDEX idx_positions_sfen_hash ON positions(sfen_hash);
\end{lstlisting}

\newpage
\section{數據科學課程報告：將棋 Policy/Value 模型}

\subsection{數據科學問題定義}
數據科學部分的任務可視為監督式學習問題。給定一個將棋局面 $s$，模型需要學習兩個目標：

\begin{enumerate}
  \item \textbf{Policy prediction：} 預測棋手在局面 $s$ 下會選擇哪一手 $a$。
  \item \textbf{Value prediction：} 預測目前輪到下棋的一方最後是否會獲勝。
\end{enumerate}

因此每筆資料可表示為：

\[
(s_i, a_i, z_i)
\]

其中 $s_i$ 是局面張量，$a_i$ 是實戰下一手的分類標籤，$z_i \in \{-1,0,1\}$ 是從目前下棋方視角定義的勝負結果。

\subsection{資料來源與資料規模}
目前專案使用 1000 份 CSA 棋譜，經資料庫匯出後形成：

\begin{longtable}{p{0.42\textwidth}p{0.42\textwidth}}
\toprule
\textbf{項目} & \textbf{數值}\\
\midrule
\endhead
原始 CSA 棋譜數 & 1000 份\\
訓練樣本數 & 111238 筆\\
局面張量形狀 & $(111238, 43, 9, 9)$\\
走法標籤數 & 13689 類\\
勝負標籤樣本數 & 111238 筆\\
訓練集 & 88473 筆\\
驗證集 & 11166 筆\\
測試集 & 11599 筆\\
\bottomrule
\end{longtable}

資料集中結果分布如下：

\begin{longtable}{p{0.42\textwidth}p{0.42\textwidth}}
\toprule
\textbf{結果類型} & \textbf{樣本數}\\
\midrule
\endhead
\texttt{TORYO} & 110967\\
\texttt{DRAW} & 104\\
\texttt{SENNICHITE} & 167\\
\bottomrule
\end{longtable}

\subsection{局面特徵工程}
將棋棋盤為 $9 \times 9$，但單純用一個整數矩陣不足以表達棋子種類、敵我關係、持駒與輪到誰下。因此本專案採用多平面 one-hot / count encoding，將每個局面編成 $43 \times 9 \times 9$ 張量。

\begin{itemize}
  \item 前 28 個平面：棋盤上雙方棋子種類。14 種棋子類型乘上己方與對方。
  \item 第 29 到 35 個平面：己方持駒數量，共 7 類可打入棋子。
  \item 第 36 到 42 個平面：對方持駒數量。
  \item 第 43 個平面：輪到哪一方下，或在旋轉視角下作為 side-to-move 輔助資訊。
\end{itemize}

設計上預設使用 \textbf{side-to-move orientation}。也就是當白方下棋時，局面與走法會旋轉 180 度，使模型永遠從「目前下棋方」角度看棋盤。這樣可以減少黑白對稱造成的資料稀疏，讓同一種戰術型態不會因方向不同而被模型視為完全不同的模式。

\subsection{走法標籤設計}
Policy head 需要將合法走法轉成固定類別。將棋走法分為一般移動與打入持駒：

\begin{itemize}
  \item 一般移動：來源格 $81$ 種、目標格 $81$ 種、是否升變 $2$ 種，共 $81 \times 81 \times 2 = 13122$ 類。
  \item 打入持駒：7 種可打入棋子、目標格 $81$ 種，共 $7 \times 81 = 567$ 類。
\end{itemize}

總標籤數為：

\[
81 \times 81 \times 2 + 7 \times 81 = 13689
\]

雖然在任一局面中合法手遠少於 13689 類，但固定輸出空間使神經網路可用標準 cross entropy loss 訓練。

\subsection{Value 標籤設計}
Value label 從目前下棋方視角定義：

\[
z =
\begin{cases}
1, & \text{目前下棋方最後獲勝}\\
0, & \text{和棋}\\
-1, & \text{目前下棋方最後落敗}
\end{cases}
\]

CSA 中的 \texttt{TORYO} 表示投了，實作上會根據最後局面的 \texttt{side\_to\_move} 推定勝方；\texttt{DRAW} 與 \texttt{SENNICHITE} 則標記為 0。後來加入的 \texttt{value\_masks} 用來表示哪些樣本有可靠 value 標籤，避免未來遇到未知結果時強迫模型學習錯誤標籤。目前資料集中 111238 筆樣本皆有 value 標籤。

\subsection{模型架構}
本專案使用小型 CNN policy/value network。輸入為 $43 \times 9 \times 9$ 張量，先經過三層 convolution 抽取局部棋型特徵，再分成 policy head 與 value head。

\begin{align*}
x &\in \mathbb{R}^{43 \times 9 \times 9}\\
h &= f_{\text{conv}}(x)\\
\hat{p} &= f_{\text{policy}}(h) \in \mathbb{R}^{13689}\\
\hat{v} &= f_{\text{value}}(h) \in [-1,1]
\end{align*}

模型結構如下：

\begin{itemize}
  \item Feature extractor：\texttt{Conv2d(43,64)}、ReLU、\texttt{Conv2d(64,64)}、ReLU、\texttt{Conv2d(64,32)}、ReLU。
  \item Policy head：flatten 後接 \texttt{Linear(32*9*9,512)}、ReLU、\texttt{Linear(512,13689)}。
  \item Value head：flatten 後接 \texttt{Linear(32*9*9,128)}、ReLU、\texttt{Linear(128,1)}、Tanh。
\end{itemize}

CNN 適合本任務的原因是棋盤具有明顯空間結構。棋子的攻防、升變區、王周圍安全性與棋子效率都與局部相對位置有關，而 convolution 能用共享權重學習這些局部模式。

\subsection{損失函數與訓練方法}
訓練時同時最佳化 policy loss 與 value loss：

\[
\mathcal{L}
=
\mathcal{L}_{policy}
+
\lambda \mathcal{L}_{value}
\]

其中：

\[
\mathcal{L}_{policy}
=
-\log
\frac{\exp(\hat{p}_{a})}
{\sum_j \exp(\hat{p}_{j})}
\]

\[
\mathcal{L}_{value}
=
m \cdot (\hat{v} - z)^2
\]

$a$ 是實戰走法標籤，$z$ 是勝負標籤，$m$ 是 value mask，$\lambda$ 為 \texttt{value-loss-weight}，目前設定為 0.25。

訓練策略包含：

\begin{itemize}
  \item 使用 AdamW optimizer，降低 overfitting 並加入 weight decay。
  \item 使用 cosine annealing learning rate scheduler，使學習率逐步下降。
  \item 以棋局為單位切分 train/validation/test，避免同一盤棋的相鄰局面同時出現在訓練與測試中造成資料洩漏。
  \item 使用 Top-1、Top-3、Top-5 accuracy 評估 policy 預測，使用 MSE 與 MAE 評估 value 預測。
\end{itemize}

\subsection{目前訓練結果}
目前模型訓練 12 epochs，checkpoint 中保存的測試集結果如下：

\begin{longtable}{p{0.42\textwidth}p{0.42\textwidth}}
\toprule
\textbf{指標} & \textbf{數值}\\
\midrule
\endhead
Test policy loss & 4.4340\\
Top-1 accuracy & 0.2230\\
Top-3 accuracy & 0.3695\\
Top-5 accuracy & 0.4385\\
Value loss & 1.0876\\
Value MAE & 0.9263\\
\bottomrule
\end{longtable}

Top-1 accuracy 約 22.3\% 代表模型直接猜中實戰下一手的比例。由於將棋合法手分支大、同一局面可能存在多個合理手，因此 Top-3 與 Top-5 更能反映 policy 是否把合理候選手排在前面。目前 Top-5 約 43.8\%，可作為 AI 搜尋的 move ordering 輔助。

\subsection{傳統搜尋與神經網路整合}
本專案不是只使用神經網路直接下棋，而是將 neural policy/value 與傳統 alpha-beta 搜尋結合。

\subsubsection{局面評估}
\texttt{ai\_search.py} 中設計了手工評估函數，包含：

\begin{itemize}
  \item 棋子價值：步、香、桂、銀、金、角、飛、王與升變棋子。
  \item 持駒價值：將棋特有的打入資源。
  \item 棋子位置價值：升變區、中央控制、棋子前進程度。
  \item 王安全：王周圍守備、敵方攻擊、線性壓力。
  \item 機動性：可走步數與大子活動能力。
  \item 懸空棋與棋子效率：被攻擊但缺乏保護的棋子會扣分。
\end{itemize}

\subsubsection{Negamax Alpha-Beta}
搜尋採用 negamax 表示法，因為雙人零和遊戲中，對手的最佳分數可視為自己的負分：

\[
V(s) = \max_{a \in A(s)} -V(T(s,a))
\]

alpha-beta pruning 使用上下界 $\alpha$ 與 $\beta$ 剪去不可能影響最終選擇的分支。此方法能在相同時間內搜尋更深層局面。

\subsubsection{搜尋優化}
專案中實作了多種搜尋優化：

\begin{itemize}
  \item \textbf{Transposition table：} 使用 zobrist hash 保存已搜尋局面，避免重複計算。
  \item \textbf{Quiescence search：} 在深度到達時繼續搜尋吃子或升變等戰術手，降低 horizon effect。
  \item \textbf{Killer move heuristic：} 保存曾造成 beta cutoff 的安靜手，提升同層搜尋排序。
  \item \textbf{History heuristic：} 對過去有效的走法累積分數，用於排序。
  \item \textbf{Policy move ordering：} 使用神經網路 policy 將合法手排序，讓 alpha-beta 優先搜尋模型認為較合理的手。
  \item \textbf{Value head bonus：} 可在根節點加入 value head 評估，補強手工評估函數。
\end{itemize}

\subsection{數據科學流程評估}
本專案的數據科學設計有三個重要特點：

\begin{enumerate}
  \item \textbf{資料不是直接從檔案訓練，而是從資料庫匯出。} 這讓資料清洗、查詢、過濾與統計更容易追蹤。
  \item \textbf{模型同時學 policy 與 value。} Policy 用於模仿棋手下一手，value 用於估計局面好壞，兩者可同時服務搜尋。
  \item \textbf{資料切分以棋局為單位。} 避免相鄰局面洩漏到測試集，使評估更接近真實泛化能力。
\end{enumerate}

目前限制則包括：

\begin{itemize}
  \item 資料量仍相對有限，1000 盤棋對深度學習將棋模型而言只適合作為小型原型。
  \item Value label 只由終局勝負回推，沒有使用中盤局面的人類評估或引擎分數，因此 value head 學習較困難。
  \item 目前 CNN 架構較小，尚未使用 residual blocks、attention 或合法手 mask loss。
  \item Policy loss 仍對所有 13689 類輸出做分類，推論時才在合法手內排序；未來可加入合法手遮罩改善學習效率。
\end{itemize}

\section{兩門課程的分工總結}

\subsection{資料庫課程重點}
資料庫課程的貢獻在於把 CSA 棋譜從文字檔轉換為關聯式資料庫。重點包含 ER model、正規化、外鍵關聯、交易控制、參數化查詢、資料去重與 API 查詢。此部分強調「如何設計資料結構，使棋譜可以被可靠保存與查詢」。

\subsection{數據科學課程重點}
數據科學課程的貢獻在於把資料庫中的棋譜轉換為機器學習樣本。重點包含特徵工程、走法標籤設計、勝負標籤、CNN policy/value network、訓練/驗證/測試切分、評估指標，以及與 alpha-beta 搜尋的整合。此部分強調「如何從結構化資料中學習決策模型」。

\subsection{共同成果}
兩門課雖然重點不同，但在本專案中互相依賴。資料庫提供乾淨、可追蹤、可查詢的資料來源；數據科學模型則將資料轉化為可用於 AI 對弈的決策能力。最終網頁系統把兩者整合成可互動展示的應用：使用者可以查棋譜、看局面、搜尋棋局、自己下棋，也可以與 AI 對弈並觀察模型候選手。

\section{結論}
shogiAI 專案展示了一個從資料工程到人工智慧應用的完整流程。資料庫部分解決了棋譜保存、查詢與重建問題；數據科學部分則將這些資料轉換成神經網路可學習的 policy/value 任務；主系統再透過 API 與前端把資料與模型串接成可操作的產品。此專案的價值不只在於完成一個將棋 AI，而是在於建立一套可擴充的資料與模型基礎。未來若加入更多棋譜、更強的模型架構、合法手遮罩、引擎評分資料與更完整的資料庫索引，系統可以進一步提升棋力、查詢效能與分析能力。

\appendix
\section{附錄：主要執行指令}

\subsection{匯入 CSA 至 MySQL}
\begin{lstlisting}[style=cmd]
python src\import_csa_to_db.py ^
  --input data ^
  --recursive ^
  --create-tables ^
  --host 140.135.65.53 ^
  --port 3306 ^
  --user 11211213 ^
  --database DB11211213 ^
  --skip-existing
\end{lstlisting}

\subsection{由 MySQL 匯出訓練資料}
\begin{lstlisting}[style=cmd]
python src\export_policy_dataset.py ^
  --output out\policy_dataset.npz ^
  --host 140.135.65.53 ^
  --port 3306 ^
  --user 11211213 ^
  --password <password> ^
  --database DB11211213
\end{lstlisting}

\subsection{訓練 Policy/Value Network}
\begin{lstlisting}[style=cmd]
python src\train_policy.py ^
  --input out\policy_dataset.npz ^
  --output out\policy_model.pt ^
  --batch-size 256 ^
  --value-loss-weight 0.25
\end{lstlisting}

\subsection{啟動棋譜瀏覽與 AI 對弈 API}
\begin{lstlisting}[style=cmd]
python src\csa_browser_api.py ^
  --host 127.0.0.1 ^
  --port 8000 ^
  --source mysql ^
  --db-host 140.135.65.53 ^
  --db-port 3306 ^
  --db-user 11211213 ^
  --db-name DB11211213 ^
  --policy-model out\policy_model.pt ^
  --value-weight 200 ^
  --policy-order-ply 1
\end{lstlisting}

\end{document}
